\section{Introduction}\label{sec:introduction}
A scheduler ranking is a conditional statement. It compares a fixed set of
implementations under a workload, an operating region, an objective and an
execution model. A ranking obtained from one trace cannot automatically be
used to select a policy for another trace; nor does agreement on average
throughput imply agreement on deadline attainment or tail latency. This
paper studies how such comparative conclusions transfer, rather than
introducing a new scheduler.

LLM serving makes this question concrete. Continuous batching and paged
key--value (KV) memory management permit concurrent requests
\citep{kwon2023vllm,yu2022orca}. Chunked prefill changes competition between
prompt processing and decoding \citep{agrawal2024sarathi}; disaggregation
separates their resource requirements \citep{zhong2024distserve,patel2024splitwise}.
Production traces differ in arrival bursts, prompt lengths and output lengths
\citep{burstgpt2025,stojkovic2025dynamollm}. These established observations
motivate testing the stability of a comparison across conditions. They do
not establish which policies will reverse order, or even that reversals
must occur in a particular evaluation.

We study the LLM-Serving Scheduler Portability Benchmark (LSSP), a frozen,
trace-driven simulator campaign. Its primary corpus contains 120 windows
from three separately released trace sources. A 13-policy execution panel,
including 11 primary policies and two style approximations, is evaluated
in six FIFO-referenced operating regions. There are 9,360 distinct
configurations and two deterministic verification passes per configuration.
The passes check repeatability of the simulator; they are not independent
samples. Historical resampling variation is instead assessed across workload windows, subject
to the dependence limitations in Section~\ref{sec:threats}.

The contribution is a reproducible evaluation methodology and an empirical
characterization of \emph{simulated} scheduler-ranking portability. The
methodology combines a fixed comparison panel, matched policy outcomes,
explicit metric eligibility, rank agreement, practical-margin reversal
classification and a finite-corpus ranking-recovery analysis. Kendall
correlation, bootstrap intervals and false-discovery correction are
standard statistical tools, not a new algorithm or theoretical result.
The scientific object is the reliability and scope of the comparative
conclusion produced by their application.

\subsection{Why comparative validity is the evaluation target}
A performance evaluator often has two decisions to make: which
implementation to compare, and which evidence would justify carrying its
ordering into a new setting. The second decision can fail even when the
first experiment is reproducible. An average over a trace can hide a
policy pair whose ordering changes elsewhere. Conversely, a small rank
change can be irrelevant if both implementations attain almost identical
utility. We therefore separate agreement of the full panel, individual
pair reversals, and changes in the set of preferred policies. These are
complementary questions with different denominators.

This distinction matters for modeling as well as for deployment. A
simulator can be repeatable while representing service demands or resource
constraints differently for different policies. Its ranking is then a
property of the complete policy--model combination. Reporting only a
leaderboard would encourage readers to attribute all differences to the
scheduling discipline. Reporting the model boundary, absolute outcome
levels and comparison support makes that attribution testable. This paper
uses the historical corpus to expose where the evidence ends rather than
treating a policy's implementation name as validation.

The unit of observation is equally consequential. Thousands of execution
records need not provide thousands of independent observations: policies
reuse the same windows, regions are transformations of those windows, and
the second deterministic pass supplies no sampling variation. Similarly,
more metric pairs do not create new workload coverage. A portable
comparison needs a clear account of what was resampled and what remained
fixed. This motivates preserving the full comparison grid, including
stable, tied and undefined cases, rather than reporting only reversals.

\subsection{Contribution and evidential scope}
The primary ANWG ranking correlations range from 0.551 to 1.000 across 18
source-pair/region comparisons. Of 990 pairwise-policy comparison records,
36 meet the saved practical-margin and nominal interval reversal criteria;
these span five policy pairs and all involve BurstGPT. Most comparisons
therefore retain their order. A separately labeled post-campaign analysis
shows a median cross-metric Kendall correlation of 0.419 and agreement of
the deterministic top-choice rule in 316 of 990 conditions. A retrospective audit additionally distinguishes deterministic choice
from the existence of a common maximizing policy, and saved interval
support from calibrated inference. Agreement and disagreement are both
findings; neither is assumed in advance.

The resulting contribution has three parts: a fixed evidence organization
that separates ranking agreement from practical pairwise change; a fully
identified simulator case study spanning source, operating region and
objective; and an independently reproducible audit showing how ties,
missingness and inferential assumptions alter the meaning of the saved
summaries. The audit is post hoc. It strengthens traceability, not the
prospective status of the experiment. No new scheduler, rank statistic,
coverage theorem or learned transfer predictor is proposed.

These are observations of historical simulator implementations. Later
source and engineering audits identified limits in partial-prefill
execution, native admission, KV allocation, preemption and execution
acknowledgment. The historical GPU comparison also has metadata and replay
limitations. It cannot establish faithful native-policy ranking transfer.
Accordingly, this paper neither claims physical validation of its simulated
reversals nor generalizes them to GPU architectures or parallel deployments.
The distinction is central to the domain of validity, not a qualification
added after interpreting the numerical results.

The four retrospective manuscript questions and their mapping to all six
original protocol questions appear in
Section~\ref{sec:rq}. Sections~\ref{sec:execution-model} and
\ref{sec:workloads} define the model and experimental design;
Sections~\ref{sec:within}, \ref{sec:results} and
\ref{sec:results-cross-metric} separate operating-region behavior,
cross-source portability and cross-metric portability. Robustness and
practical implications precede an explicit account of threats to validity.
